\section{引言与 Vision Transformer 回顾}

本节课覆盖计算机视觉中若干核心任务——语义分割、目标检测、实例分割，以及神经网络的可视化与理解方法。在正式进入这些主题之前，讲者首先简要回顾了上节课介绍的 Vision Transformer（ViT）架构。

\begin{figure}[H]
\centering
\includegraphics[width=0.95\textwidth]{slides-images/slide-001.jpg}
\caption{课程大纲：本节涵盖目标检测、语义分割、实例分割和可视化理解\protect\footnotemark}
\end{figure}
\footnotetext{Slides 第2页。}

\subsection{ViT 架构核心}

Vision Transformer 的核心思想是将图像分割为若干 patch（例如 $S \times S$ 的网格），将每个 patch 通过线性投影映射为 $D$ 维的 token，加上位置编码（positional embedding）后送入标准 Transformer encoder。

\begin{figure}[H]
\centering
\includegraphics[width=0.95\textwidth]{slides-images/slide-003.jpg}
\caption{ViT 架构：图像被切分为 patch，线性投影为 token，加上位置编码后送入 Transformer 层\protect\footnotemark}
\end{figure}
\footnotetext{Slides 第4页。}

以一张 $224 \times 224 \times 3$ 的输入图像为例，ViT 的具体数据流如下：

\begin{enumerate}
    \item 将图像切分为 $14 \times 14 = 196$ 个大小为 $16 \times 16 \times 3$ 的 patch
    \item 每个 patch 展平后为 $768$ 维向量，通过线性变换映射为 $D$ 维 token
    \item 加上可学习的位置编码（$D$ 维向量），告知 Transformer 每个 patch 的二维空间位置
    \item 将 196 个 token 送入标准 Transformer encoder（自注意力 + MLP）
\end{enumerate}

\begin{figure}[H]
\centering
\includegraphics[width=0.95\textwidth]{slides-images/slide-016.jpg}
\caption{ViT 的详细流程：$224 \times 224 \times 3$ 图像切分为 $16 \times 16$ patch，线性投影到 $D$ 维，再经过 Transformer 层处理\protect\footnotemark}
\end{figure}
\footnotetext{Slides 第17页。}

\begin{importantbox}{ViT 分类的两种方式}
\begin{enumerate}
    \item \textbf{Class Token}：添加一个可学习的特殊 token 作为输入（$D$ 维），其对应的输出经过线性层映射为类别概率向量
    \item \textbf{全局池化}：对所有输出 token 做 pooling，再映射为类别概率
\end{enumerate}
两种方式均使用 softmax 交叉熵损失进行监督训练。关键点在于：ViT 不使用任何 masking——每个 patch 可以关注所有其他 patch，这与 NLP 中的 causal attention 不同。
\end{importantbox}

\begin{knowledgebox}{ViT 与 CNN 的关键差异}
\begin{itemize}
    \item \textbf{感受野}：CNN 的感受野从局部逐步扩大；ViT 从第一层就具有全局感受野
    \item \textbf{归纳偏置}：CNN 具有平移不变性和局部连接假设；ViT 几乎没有这些先验
    \item \textbf{数据需求}：正因缺少归纳偏置，ViT 通常需要更大规模的数据集才能充分训练
    \item \textbf{可扩展性}：ViT 的自注意力复杂度为 $O(N^2)$（$N$ 为 token 数），而 CNN 为 $O(N)$
\end{itemize}
\end{knowledgebox}

\subsection{现代 Transformer 的优化技巧}

讲者介绍了几种在现代 Transformer 架构中广泛使用的优化手段：

\begin{figure}[H]
\centering
\includegraphics[width=0.95\textwidth]{slides-images/slide-008.jpg}
\caption{Transformer 优化技巧汇总：Pre-Norm、RMS Norm、SwiGLU MLP、Mixture of Experts\protect\footnotemark}
\end{figure}
\footnotetext{Slides 第9页。}

\textbf{1. Pre-Norm（前置归一化）}：将 Layer Norm 放在自注意力和 MLP 之前（而非之后），以保留残差连接的恒等映射能力，使训练更加稳定。原始 Transformer 使用 Post-Norm（先计算再归一化），现代模型几乎全部改用 Pre-Norm。

\textbf{2. RMS Norm}：使用均方根归一化代替标准 Layer Norm，省去均值中心化步骤：

$$y_i = \frac{x_i}{\text{RMS}(x)} \cdot \gamma_i, \quad \text{RMS}(x) = \sqrt{\varepsilon + \frac{1}{N}\sum_{i=1}^{N} x_i^2}$$

\begin{itemize}
    \item $x$：输入向量（shape $D$）
    \item $\gamma$：可学习的缩放参数（shape $D$）
    \item $\varepsilon$：防止除零的小常数
\end{itemize}

\begin{figure}[H]
\centering
\includegraphics[width=0.95\textwidth]{slides-images/slide-022.jpg}
\caption{RMS Norm 的公式与在 Transformer 中的位置：替换 Layer Norm，训练更稳定\protect\footnotemark}
\end{figure}
\footnotetext{Slides 第23页。}

\textbf{3. SwiGLU MLP}：使用门控非线性（gated nonlinearity），引入第三个权重矩阵 $W_3$，在保持参数量不变的情况下（隐藏层维度设为 $8d/3$）学习更高维的非线性表示。

\textbf{4. Mixture of Experts (MoE)}：将 MLP 层替换为多组并行的 MLP ``专家''，通过 router 将不同 token 路由到不同专家。这增加了参数量但不成比例地增加计算量（每次只有部分专家被激活），被现代大型语言模型广泛采用。

\begin{table}[H]
\centering
\begin{tabular}{p{3cm}p{4.5cm}p{5cm}}
\toprule
\textbf{技巧} & \textbf{核心改进} & \textbf{典型效果} \\
\midrule
Pre-Norm & 归一化前置于残差分支 & 训练更稳定，不需要 warm-up \\
RMS Norm & 去掉均值归一化步骤 & 计算更快，精度不降 \\
SwiGLU MLP & 门控非线性，三矩阵设计 & 同参数量下表达力更强 \\
MoE & 多专家并行，稀疏激活 & 参数量大增但计算量可控 \\
\bottomrule
\end{tabular}
\caption{现代 Transformer 四大优化技巧对比}
\end{table}

\begin{warningbox}{MoE 并非免费午餐}
MoE 虽然在理论上可以用固定计算量支撑更大的参数空间，但实际部署中面临几个挑战：(1) 专家负载不均衡可能导致部分专家过载；(2) 模型的总参数量更大，推理时需要加载所有专家权重到内存；(3) 路由策略本身需要精心调优。
\end{warningbox}

\subsection{本章小结}

Vision Transformer 将图像视为 patch 序列，通过自注意力机制全局建模像素间关系。以 $224 \times 224$ 的图像为例，ViT 会生成 196 个 token，每个 token 从第一层起就具有全局感受野。现代 Transformer 通过 Pre-Norm、RMS Norm、SwiGLU 和 MoE 等技巧进一步提升了训练稳定性和模型容量，这些改进构成了当前视觉和语言模型的标准组件。

%% ========== Part 2: 语义分割 ==========

